\section{Decision Framework}
\label{sec:decision}

\subsection{Overview}

G.14 provided a four-question decision tree for selecting the search algorithm.
This section extends the framework to six questions covering both structure and
search layers.
Questions Q1--Q2 select the structure algorithm; questions Q3--Q6 select the
search algorithm (matching the G.14 four-question tree, with one modification
at Q3 to incorporate \pt{} and \ams{}).

The framework encodes structural properties of the algorithms established in
their respective B/G/H-series papers, not solely the single-run empirics of
this paper.

\subsection{Structure Selection: Q1--Q2}

\begin{enumerate}[label=\textbf{Q\arabic*.}]

  \item \textbf{Is certified-optimal EC required?}

        If the redistricting task requires a provably minimum-\ec{} initial
        plan (e.g., for small-state legal challenges where optimality is
        litigated, or for academic benchmarking):

        \begin{itemize}
          \item If $N \leq 500$ per bisection half (applicable to NH $k = 2$
                and very small custom jurisdictions): use \ilp{} (B.24).
          \item If $N > 500$: ILP is not tractable. Use SA (B.19) for the
                best available heuristic EC; the optimality gap is
                approximately 8--12\% relative to ILP.
        \end{itemize}

        \textit{Note:} edge-cut is not the standard compactness metric in
        U.S. redistricting law; courts more commonly apply Polsby-Popper,
        Reock, or visual compactness.
        The ILP certificate (``no valid plan achieves lower edge cuts'') is
        useful for academic benchmarking and may support expert testimony
        in jurisdictions where edge-cut compactness is litigated, but does
        not directly satisfy Polsby-Popper or visual-compactness legal tests.
        Consult B.24~\S5 for the precise certificate scope.

        \emph{Proceed to: \ilp{} or \sa{}. Then proceed to Q3 for search.}

        If certified or near-optimal EC is not required, proceed to Q2.

  \item \textbf{Is geometric compactness (PP) the primary structure
        criterion?}

        If the redistricting standard references ``compactness'' in the
        geometric or perimeter-area sense (as in most legal standards citing
        Polsby-Popper), the structure algorithm is the primary PP determinant:
        search algorithms raise PP modestly from the structure baseline
        (typically $< 8\%$ improvement).

        \begin{itemize}
          \item If centroid data is available (standard for 2020 decennial
                pipeline): use \cvdgeo{} (B.22b).
          \item If centroid data is unavailable or the input graph is
                non-standard: use \cvdg{} (B.22).
          \item If PP is important but runtime is constrained: use
                \bfsgrow{} (B.23); competitive PP, fastest structure method.
        \end{itemize}

        \emph{Proceed to: \cvdgeo{}, \cvdg{}, or \bfsgrow{}.
        Then proceed to Q3 for search.}

        If neither certified EC nor elevated PP is a primary concern,
        use \metis{} (default; fastest, lowest structure overhead)
        and proceed to Q3.

\end{enumerate}

\subsection{Search Selection: Q3--Q6}

\begin{enumerate}[label=\textbf{Q\arabic*.},start=3]

  \item \textbf{Is VRA compliance required during sampling?}

        If minority-majority district preservation must be enforced during
        chain execution (not merely verified post hoc):
        \emph{Proceed to: \vra{} (G.13). Stop.}

        If VRA is verified post hoc, proceed to Q4.

  \item \textbf{Is $k \geq 15$?}

        For large $k$ ($\geq 15$: TX, FL, CA, NY, PA, OH, and others):

        \begin{itemize}
          \item If \textbf{compactness optimisation} (lowest EC) is the goal:
                use \pt{} (B.20). Requires multi-core hardware (8--16 cores).
                For single-core constraint: use \sburst{} (G.6).
          \item If \textbf{fast mixing} (low ACF, broad distributional
                coverage) is the goal: use \ams{} (B.21) for self-tuning, or
                \ms{} (G.11) for fixed $\alpha$.
          \item If \textbf{calibrated distribution} is required: use \smc{}
                (G.7) or \smcp{} (new).
        \end{itemize}

        For $k < 15$, single-scale chains mix adequately; proceed to Q5.

  \item \textbf{Is a calibrated target distribution required?}

        \begin{itemize}
          \item Exact calibration (importance-weighted, user-specified):
                use \smcp{} for a percentile-selected plan, or \smc{} for
                an arbitrary plan from the particle cloud.
                Difference: \smcp{} returns the median plan (or $p$-th
                percentile); \smc{} returns the final particle.
                For court submissions, either is defensible; \smcp{} is
                preferred when the point estimate must be described precisely.
          \item Approximate calibration (MH-corrected chain): use \msp{}
                (G.10) as the default. \fr{} (G.9) is the alternative.
        \end{itemize}

        If distributional calibration is not required, proceed to Q6.

  \item \textbf{Is pure compactness optimisation the goal?}

        \begin{itemize}
          \item Best EC, no distributional constraint: use \sburst{} (G.6).
          \item Best EC with MH-corrected distributional properties: use
                \sburstf{} (G.12).
        \end{itemize}

        If the goal is not compactness optimisation, the remaining algorithm
        is \flip{} (G.8): local boundary sensitivity analysis.
        \emph{Proceed to: \flip{} (G.8). Stop.}

\end{enumerate}

\subsection{Decision Tree Diagram}

\begin{center}
\begin{tikzpicture}[
  decision/.style={
    diamond, draw, fill=blue!8, text width=3.6cm, align=center,
    inner sep=1pt, font=\small
  },
  result/.style={
    rectangle, draw, fill=green!10, text width=3.2cm, align=center,
    rounded corners=3pt, font=\small
  },
  arrow/.style={->, >=Stealth, thick},
  label/.style={font=\small\itshape},
]

%% Structure layer header
\node[font=\bfseries\small, text=blue!60!black] (shdr) at (0, 0)
  {Structure Layer};

%% Q1
\node[decision, below=0.5cm of shdr] (q1)
  {Certified optimal EC?\\($N \le 500$?)};

\node[result, right=2.0cm of q1, yshift=1.2cm] (ilp)
  {\ilp{} (B.24)\\\footnotesize Exact optimality};
\draw[arrow] (q1.north east)
  -- node[label, above, sloped]{\small Yes} (ilp.west);

\node[result, right=2.0cm of q1, yshift=-0.5cm] (sa)
  {\sa{} (B.19)\\\footnotesize Best heuristic EC};
\coordinate (branch1) at ($(q1.east) + (0.8cm, 0)$);
\draw[arrow] (q1.east) -- (branch1);
\draw[arrow] (branch1) |- node[label, right, pos=0.75]{\small N too large} (sa.west);

%% Q2
\node[decision, below=1.5cm of q1] (q2) {PP primary?\\Centroids available?};
\draw[arrow] (q1.south) -- node[label, right]{\small No} (q2.north);

\node[result, right=2.0cm of q2, yshift=1.4cm] (cvdgeo)
  {\cvdgeo{} (B.22b)\\\footnotesize Best PP};
\node[result, right=2.0cm of q2, yshift=-0.2cm] (bfs)
  {\bfsgrow{} (B.23)\\\footnotesize Fast, good PP};
\coordinate (branch2) at ($(q2.east) + (0.8cm, 0)$);
\draw[arrow] (q2.east) -- node[label, above]{\small Yes} (branch2);
\draw[arrow] (branch2) |- node[label, right, pos=0.75]{\small Centroids} (cvdgeo.west);
\draw[arrow] (branch2) |- node[label, right, pos=0.75]{\small Fast} (bfs.west);

\node[result, below=1.0cm of q2, xshift=2.5cm] (metis)
  {\metis{} (default)\\\footnotesize Fastest};
\draw[arrow] (q2.south) -- node[label, right]{\small No} (metis.north);

%% Search layer header
\node[font=\bfseries\small, text=red!60!black, below=1.8cm of metis] (srhdr)
  {Search Layer};

%% Q3
\node[decision, below=0.5cm of srhdr] (q3) {VRA required\\during sampling?};

\node[result, right=2.0cm of q3, yshift=1.0cm] (vra)
  {\vra{} (G.13)\\\footnotesize VRA-constrained};
\draw[arrow] (q3.north east)
  -- node[label, above, sloped]{\small Yes} (vra.west);

%% Q4
\node[decision, below=1.5cm of q3] (q4) {$k \geq 15$?};
\draw[arrow] (q3.south) -- node[label, right]{\small No} (q4.north);

\node[result, right=2.0cm of q4, yshift=1.4cm] (pt)
  {\pt{} / \ms{} (B.20/G.11)\\\footnotesize Large-$k$ mixing};
\draw[arrow] (q4.north east)
  -- node[label, above, sloped]{\small Yes} (pt.west);

%% Q5
\node[decision, below=1.5cm of q4] (q5) {Calibrated\\distribution?};
\draw[arrow] (q4.south) -- node[label, right]{\small No} (q5.north);

\node[result, right=2.0cm of q5, yshift=1.4cm] (smcp2)
  {\smcp{} / \smc{} (G.7)\\\footnotesize Exact calibration};
\node[result, right=2.0cm of q5, yshift=-0.4cm] (ms2)
  {\msp{}/\fr{} (G.10/9)\\\footnotesize Approx.\ uniform};
\coordinate (branch5) at ($(q5.east) + (0.8cm, 0)$);
\draw[arrow] (q5.east) -- node[label, above]{\small Yes} (branch5);
\draw[arrow] (branch5) |- node[label, right, pos=0.75]{\small Exact} (smcp2.west);
\draw[arrow] (branch5) |- node[label, right, pos=0.75]{\small Approx} (ms2.west);

%% Q6
\node[decision, below=1.5cm of q5] (q6) {Compactness\\optimisation?};
\draw[arrow] (q5.south) -- node[label, right]{\small No} (q6.north);

\node[result, right=2.0cm of q6, yshift=1.0cm] (sb)
  {\sburst{}/\sburstf{} (G.6/12)\\\footnotesize Compact plans};
\draw[arrow] (q6.north east)
  -- node[label, above, sloped]{\small Yes} (sb.west);

\node[result, below=1.0cm of q6] (flipr)
  {\flip{} (G.8)\\\footnotesize Sensitivity};
\draw[arrow] (q6.south) -- node[label, right]{\small No} (flipr.north);

\end{tikzpicture}
\end{center}

\subsection{Combined Structure + Search Recommendation Table}

Table~\ref{tab:combined} gives the recommended structure + search combination
for each practitioner use case.

\begin{table}[ht]
\centering
\caption{%
  Combined structure + search recommendations by use case.
  ``Alternative'' is the second choice when the primary is unavailable.
  All flags are for the \texttt{redist} CLI.
}
\label{tab:combined}
\small
\begin{tabular}{@{}p{4.5cm}p{4.2cm}p{4.2cm}@{}}
\toprule
Use Case & Recommended & Alternative \\
\midrule
Court submission\\(calibrated, neutral)
  & \texttt{--structure prime-factor}\newline
    \texttt{--search smc-percentile}
  & \texttt{--structure prime-factor}\newline
    \texttt{--search merge-split} \\[6pt]
EC minimisation,\\small state ($k \leq 4$, $N \leq 500$)
  & \texttt{--structure ilp}\newline
    \texttt{--search single}
  & \texttt{--structure sa-bisect}\newline
    \texttt{--search short-burst} \\[6pt]
EC minimisation,\\large state
  & \texttt{--structure sa-bisect}\newline
    \texttt{--search parallel-tempering}
  & \texttt{--structure sa-bisect}\newline
    \texttt{--search short-burst} \\[6pt]
PP maximisation
  & \texttt{--structure cvd-geographic}\newline
    \texttt{--search smc-percentile}
  & \texttt{--structure bfs-growth}\newline
    \texttt{--search merge-split} \\[6pt]
Large-$k$ ($k \geq 15$),\\fast mixing
  & \texttt{--structure prime-factor}\newline
    \texttt{--search adaptive-multiscale}
  & \texttt{--structure prime-factor}\newline
    \texttt{--search multiscale} \\[6pt]
Large-$k$, EC optimisation\\(multi-core available)
  & \texttt{--structure sa-bisect}\newline
    \texttt{--search parallel-tempering}
  & \texttt{--structure standard-bisect}\newline
    \texttt{--search short-burst} \\[6pt]
VRA Section~2 analysis
  & \texttt{--structure ratio-optimal}\newline
    \texttt{--weights vra-aligned}\newline
    \texttt{--search vra-recom}
  & --- \\[6pt]
Local sensitivity\\at a fixed plan
  & \texttt{--structure ratio-optimal}\newline
    \texttt{--search flip}
  & --- \\[6pt]
Production sweep,\\general purpose
  & \texttt{--structure prime-factor}\newline
    \texttt{--weights county}\newline
    \texttt{--search convergence}
  & \texttt{--structure prime-factor}\newline
    \texttt{--search merge-split} \\[6pt]
Fastest result\\(time budget limited)
  & \texttt{--structure bfs-growth}\newline
    \texttt{--search convergence}
  & \texttt{--structure standard-bisect}\newline
    \texttt{--search single} \\
\bottomrule
\end{tabular}
\end{table}

\subsection{Metric Sensitivity Table}

Table~\ref{tab:sensitivity} summarises which algorithm layer and type most
affects each metric.
This table is the primary guide for practitioners whose task is defined by
a single dominant metric.

\begin{table}[ht]
\centering
\caption{%
  Algorithm sensitivity by metric. Layer: which layer (Structure, Search, or Both)
  most determines the metric. Leader: the algorithm achieving the best
  value on this metric in this comparison.
}
\label{tab:sensitivity}
\begin{tabular}{@{}llll@{}}
\toprule
Metric & Layer & Leader & Second \\
\midrule
EC (final plan)   & Both & \pt{} + \sa{} & \sburst{} + \sa{} \\
PP (final plan)   & Structure (primary) & \cvdgeo{} + \smcp{} & \cvdgeo{} + \smc{} \\
D-seat            & Search (modest)     & Varies by state & --- \\
ACF(100)          & Search only         & \ams{} & \pt{} \\
Runtime (total)   & Both & \bfsgrow{} + \flip{} & \metis{} + \flip{} \\
\bottomrule
\end{tabular}
\end{table}

\subsection{Limitations}

The framework inherits the G.14 limitations (single-run, state-specific effects,
parameter sensitivity) plus two additional caveats for the structure layer.

\paragraph{Structure and search are not independent.}
The combined recommendation table treats the two layers as separable, but there
are interactions.
SA structure + Short-Burst search may not compound: SA already reduces EC at
structure time, and Short-Burst's burst strategy evaluates relative EC
improvement from the starting plan.
If the starting plan is already near the compact region, burst improvement
may be smaller than from a METIS baseline.
This interaction has not been formally measured; it is an open empirical
question for future work.

\paragraph{ILP intractability boundary.}
The $N \leq 500$ threshold for ILP tractability is empirical (runtime
$< 60$ s in our experiments).
For graphs with unusual structure (dense connectivity, deep tree branching),
ILP may time out at smaller $N$.
For production use, the \texttt{redist} ILP implementation includes a timeout
flag (\texttt{--ilp-timeout}) that falls back to SA when exceeded.

\paragraph{Future algorithms.}
The framework covers all 15 algorithms implemented as of May 2026.
Future algorithms will require updates; the modular six-question structure
is designed to accommodate new entries without restructuring the existing
decision logic.
